% ============================================================
%  ArcLoom SPPU — Competitive Landscape Analysis
%  Version 1.0 — May 2026
%  Confidential — Internal Use Only (DARPA/Patent Preparation)
% ============================================================
\documentclass[11pt, a4paper]{article}

\usepackage[margin=2.2cm]{geometry}
\usepackage{amsmath, amssymb}
\usepackage{booktabs}
\usepackage{array}
\usepackage{longtable}
\usepackage{xcolor}
\usepackage{hyperref}
\usepackage{titlesec}
\usepackage{enumitem}
\usepackage{fancyhdr}
\usepackage{graphicx}
\usepackage{multicol}

\hypersetup{colorlinks=true, linkcolor=black, urlcolor=blue}

\definecolor{tfe-gold}{RGB}{180, 140, 60}
\definecolor{tfe-dark}{RGB}{30, 30, 30}
\definecolor{tfe-gray}{RGB}{100, 100, 100}
\definecolor{gap-red}{RGB}{180, 40, 40}
\definecolor{proven-green}{RGB}{40, 130, 40}

\pagestyle{fancy}
\fancyhf{}
\rhead{\textcolor{tfe-gray}{\small ArcLoom Competitive Landscape v1.0}}
\lhead{\textcolor{tfe-gray}{\small Confidential --- DARPA/Patent Preparation}}
\rfoot{\textcolor{tfe-gray}{\small Page \thepage}}

\titleformat{\section}{\large\bfseries\color{tfe-dark}}{}{0em}{}[\titlerule]
\titleformat{\subsection}{\normalsize\bfseries\color{tfe-dark}}{}{0em}{}
\titleformat{\subsubsection}{\normalsize\itshape\color{tfe-dark}}{}{0em}{}

\setlength{\parindent}{0pt}
\setlength{\parskip}{6pt}

\newcommand{\nopriorart}{\textcolor{gap-red}{\textbf{No prior art}}}
\newcommand{\proven}{\textcolor{proven-green}{\textbf{PROVEN}}}

\begin{document}

\begin{center}
  {\LARGE \textbf{ArcLoom SPPU}}\\[6pt]
  {\large Competitive Landscape Analysis}\\[4pt]
  {\large Ternary Computing, Clockless Perception, and Related Work}\\[8pt]
  {\color{tfe-gray}\small Version 1.0 \quad|\quad May 2026 \quad|\quad Confidential}\\[4pt]
  {\color{tfe-gray}\small Joseph Forrester \quad|\quad DSF-AI Architecture}
\end{center}

\vspace{1em}
\hrule
\vspace{1em}

% ============================================================
\section{Purpose}

This document surveys all known ternary computing systems, clockless
perception hardware, structural memory architectures, and related work.
Its purpose is to establish that ArcLoom occupies a unique position in
the landscape and to prepare clear distinctions for DARPA pitch,
Efabless submission, and patent prosecution.

All ArcLoom claims reference validated demonstrations documented in
the ASIC Architecture Specification v1.2.

% ============================================================
\section{Ternary Computing Hardware}

\subsection{Huawei Ternary Logic Gate (2025)}

\textbf{What:} Patent CN119652311A for ternary logic gates using
threshold voltage layering on CNTFET. Claims 40\% fewer transistors,
60\% less power.

\textbf{Ternary type:} Standard (0, 1, 2) --- NOT balanced.

\textbf{Hardware status:} Patent filing. Unclear if silicon exists.

\textbf{Computing model:} Sequential instruction execution with
3-state gates.

\textbf{ArcLoom distinction:} Huawei improved the transistor. The
computing model is unchanged --- fetch/decode/execute with a clock.
Standard ternary (0,1,2) lacks the symmetric negation property of
balanced ternary ($-1$, 0, $+1$). The SPPU's combinational settling
through signed coupling is architecturally impossible in Huawei's
number system. (See ASIC Spec v1.2, \S10.2.1.)

\subsection{5500FP Ternary RISC Processor (2026)}

\textbf{What:} 24-trit balanced ternary RISC processor on Efinix
Trion T120F484 FPGA. 120-instruction ISA. Open hardware.

\textbf{Ternary type:} Balanced ($-1$, 0, $+1$).

\textbf{Hardware status:} Running on FPGA. No custom silicon.

\textbf{Computing model:} Von Neumann RISC --- fetch, decode, execute,
store.

\textbf{ArcLoom distinction:} The 5500FP treats trits as numbers in
registers. ArcLoom treats them as coupled physical states in a field.
The null trit in 5500FP is arithmetic zero. In SPPU, it is structural
uncertainty --- the system's admission that coupling pressure is
insufficient to commit this dimension. The 5500FP validates that
balanced ternary works on FPGA; ArcLoom builds a different
computational model on top of it. (See ASIC Spec v1.2, \S10.2.2.)

\subsection{Microsoft BitNet b1.58 (2024--2026)}

\textbf{What:} LLM with ternary weight quantization $\{-1, 0, +1\}$.
2B parameter model. Open-source inference engine.

\textbf{Ternary type:} Ternary VALUES on binary HARDWARE.
Not ternary computing.

\textbf{Hardware status:} Runs on standard x86 CPUs. No custom hardware.

\textbf{ArcLoom distinction:} BitNet uses ternary as a compression
trick --- reducing model weight storage from 16-bit floats to 1.58-bit
trits. The processor is entirely binary. The computation is still
matrix multiplication (simplified to add/subtract). ArcLoom uses
balanced ternary as the native computational substrate. The coupling
fabric IS the computation. No matrix multiplication. No training.
No inference.

\subsection{TCN-CUTIE / CUTIE (ETH Zurich, 2022)}

\textbf{What:} Fabricated ASIC (GlobalFoundries 22nm) for ternary
neural network acceleration. 1,036 TOp/s/W efficiency.

\textbf{Ternary type:} Ternary weights, binary logic. The chip is a
binary accelerator optimized for ternary-quantized weight matrices.

\textbf{Hardware status:} Real silicon. Production-grade.

\textbf{ArcLoom distinction:} TCN-CUTIE is a binary chip that
processes ternary data. ArcLoom is a ternary chip that perceives
through coupling. TCN-CUTIE requires trained neural network weights.
ArcLoom derives weights from sensor physics with zero training.

\subsection{TerEffic (2025)}

\textbf{What:} FPGA architecture for ternary LLM inference.
16,300 tokens/second.

\textbf{ArcLoom distinction:} Ternary data on binary hardware with
a trained model. ArcLoom is native ternary logic, no training.

\subsection{Tiny Tapeout Balanced Ternary Calculator (USN Group)}

\textbf{What:} 2-trit balanced ternary calculator (add/subtract/multiply)
taped out on Skywater 130nm through Tiny Tapeout 03. Follow-up REBEL-2
ALU on TT05.

\textbf{Ternary type:} Balanced ($-1$, 0, $+1$). Binary-Encoded
Ternary on CMOS.

\textbf{Hardware status:} Real silicon (Tiny Tapeout shuttle).

\textbf{ArcLoom distinction:} REBEL-2 does add/sub/mul. MathLoom
does all of those PLUS folding division (41,430 tests, zero errors)
with native free rounding via truncation. Nobody else has verified
balanced ternary division on hardware. More fundamentally, REBEL-2
is a standalone ALU. MathLoom exists only at I/O boundaries of the
SPPU --- the perception path has no ALU at all.

\subsection{Carbon Nanotube Ternary Logic (Science Advances, Jan 2025)}

\textbf{What:} CNT source-gating transistors with native 3-state
switching. Ternary inverters, NMIN/NMAX gates, ternary SRAM, ternary
neural network.

\textbf{Hardware status:} Lab-scale fabricated devices. Not a system.

\textbf{ArcLoom distinction:} Individual ternary gates, not an
architecture. Their ternary neural network is trained. ArcLoom is a
complete perception system from sensor to motor, training-free.

\subsection{USN Ternary Research Group (Univ.\ of South-Eastern Norway)}

\textbf{What:} Academic research group building EDA tools for ternary
VLSI. Mixed Radix Circuit Synthesizer (MRCS). Founded 2019.

\textbf{Hardware status:} EDA tools and simulation. No fabricated chips.

\textbf{ArcLoom distinction:} USN builds the tools to design ternary
chips. ArcLoom IS a ternary chip (on FPGA, ASIC spec'd). Complementary,
not competing. Potential collaboration for ArcLoom-256 tape-out.

\subsection{Historical: Setun (Moscow State University, 1958)}

\textbf{What:} First and only mass-produced balanced ternary computer.
$\sim$50 units built. Discontinued for political reasons, not technical ones.

\textbf{ArcLoom distinction:} Setun was a general-purpose computer that
happened to use balanced ternary. It validated the number system. It did
not explore coupling, perception, or the null trit as structural
uncertainty.

\subsection{TCAM (Cisco, Broadcom --- Production)}

\textbf{What:} Ternary Content-Addressable Memory in network routers.
Three states: 0, 1, X (don't care). Massively deployed.

\textbf{ArcLoom distinction:} TCAM's third state is a wildcard mask
for pattern matching, not a computational state. TCAM does lookup, not
logic. The ``ternary'' in TCAM has nothing in common with balanced
ternary computing.

\subsection{Three Generations of Ternary}

\begin{center}
\begin{tabular}{@{}llll@{}}
\toprule
\textbf{Generation} & \textbf{Architecture} & \textbf{Innovation} & \textbf{Limitation} \\
\midrule
Huawei (2025) & Standard ternary gates & Better transistor & Same computing model \\
5500FP (2026) & Balanced ternary RISC & Better number system & Same computing model \\
ArcLoom (2026) & BT coupling fabric & Different computing model & --- \\
\bottomrule
\end{tabular}
\end{center}

Huawei improved the \textit{transistor}. The 5500FP improved the
\textit{number system}. ArcLoom changed what \textit{computation means}.

% ============================================================
\section{Clockless / Asynchronous Perception Hardware}

\subsection{Intel Loihi 2}

\textbf{What:} 128 asynchronous neuron cores on Intel 4 process.
Spiking neural network processor. Sensor fusion demonstrated
(radar, vision, lidar).

\textbf{ArcLoom distinction:} Loihi runs trained spiking neural
networks on async hardware. The hardware is clockless; the computation
is ML. SPPU is combinational --- no spikes, no neurons, no training,
no iteration. The answer settles through wire propagation in $\sim$5~ns.

\subsection{BrainChip Akida (AKD1500)}

\textbf{What:} Event-driven digital processor. 800 GOPS at $<$300mW.
Gesture recognition demonstrated with Prophesee event camera.

\textbf{ArcLoom distinction:} Akida runs Temporal Event-based Neural
Networks (TENNs). Event-driven hardware, ML algorithms. SPPU has no
neural network of any kind.

\subsection{SpiNNaker2 (TU Dresden)}

\textbf{What:} 153 ARM cores with async event routing. 22nm FDSOI.
Deployed at Sandia National Labs.

\textbf{ArcLoom distinction:} SpiNNaker2 is a massively parallel ARM
cluster for spiking neural network simulation. It has 153 clocked ARM
cores connected by async routing. SPPU has zero clocked cores and zero
neural networks.

\subsection{Prophesee GenX320 (Event Camera)}

\textbf{What:} 320$\times$320 event camera where each pixel fires
independently and asynchronously. 10,000 FPS equivalent. $>$140dB
dynamic range.

\textbf{ArcLoom distinction:} Prophesee is an async \textit{sensor},
not an async \textit{processor}. Interesting as a future BSIL input
source --- event-driven pixels feeding directly into the SPPU's
combinational fabric. Complementary technology.

\subsection{Gap Analysis}

Every existing clockless/async processor runs trained neural networks.
The async fabric is the delivery mechanism; the computation is still ML.

\textbf{Nobody has:} purely combinational perception where the answer
settles through coupling constraints with no spikes, no neurons, no
weights learned from data, and no iteration. The SPPU is unique in
this category.

% ============================================================
\section{Content-Addressable Structural Memory (Krimelack)}

\subsection{Existing CAM Technology}

All deployed CAMs (TCAM in routers, memristor CAMs, ferroelectric CAMs)
do bit-level or weight-level matching. They store addresses, keys, or
feature vectors. None store \textbf{structural motifs} --- coupled
multi-field patterns that represent a settled loom state.

\subsection{Memristor Habituation (Nature Communications, 2025)}

\textbf{What:} Third-order memristor (HfO$_2$+TiO$_x$) with physical
habituation. Robot arm ignores 71\% of familiar stimuli while responding
to threats. Non-volatile habituation survives power loss.

\textbf{ArcLoom distinction:} This is the closest existing work to
Krimelack's habituation mechanism. The difference: memristor habituation
is \textit{analog}, \textit{per-synapse}, and embedded in a trained SNN.
Krimelack habituation is \textit{digital}, \textit{pattern-level}
(full loom state), and operates through \textit{combinational dead-zone
feedback}. The memristor habituates individual connections. Krimelack
habituates entire structural experiences.

\subsection{Gap Analysis}

No existing hardware combines:
\begin{enumerate}[nosep]
  \item Content-addressable structural motif storage
  \item Cue-based recall from partial signatures
  \item Pattern-level habituation via dead-zone feedback
  \item Autonomous commit gating based on dimensional exhaustion
        (L6 structural lock)
\end{enumerate}

Each piece exists in isolation somewhere. The combination is unique
to Krimelack.

% ============================================================
\section{Balanced Ternary Arithmetic (MathLoom)}

\subsection{REBEL-2 (USN Group, Tiny Tapeout 05)}

\textbf{What:} Balanced ternary ALU with MIN, MAX, ADD, SUB, MUL,
CMP, SHFT. Submitted for ASIC fabrication on Skywater 130nm.

\textbf{MathLoom distinction:} REBEL-2 does not implement division.
MathLoom's folding division --- verified with 41,430 tests (integer,
complex, matrix, polynomial) at zero errors --- has no equivalent in
any balanced ternary hardware. The native free rounding property
(truncation = rounding in balanced ternary) is exploited as a design
feature in MathLoom but is not used in REBEL-2's operations.

\subsection{Georgia Tech Discrete ALU (Louis Duret-Robert)}

\textbf{What:} 2-trit balanced ternary ALU from discrete CD4007 CMOS
chips.

\textbf{MathLoom distinction:} Proof of concept on discrete components.
Not integrated. No division. No verification suite.

\subsection{MathLoom Unique Claims}

\begin{enumerate}[nosep]
  \item Folding division with native free rounding --- \nopriorart
  \item Exact $1/3$ representation exploited as design feature --- \nopriorart
  \item Computational completeness proven (complex, matrix, polynomial, trig)
  \item 41,430 verification tests + 18 silicon tests at zero errors
  \item Positioned as I/O boundary arithmetic, not central ALU
\end{enumerate}

% ============================================================
\section{Structural Perception Without ML}

\subsection{Hyperdimensional Computing (UC San Diego, UC Berkeley, ETH Zurich)}

\textbf{What:} Encodes sensor data into high-dimensional binary vectors.
Does similarity matching via dot products. HyperSense (2024) does
real-time raw sensor processing on FPGA. GraspHD (2024) does robotic
grasping.

\textbf{ArcLoom distinction:} HDC encodes deterministically but
classifies with a trained model. The encoding is training-free; the
decision is not. SPPU is training-free end-to-end --- from sensor to
motor. HDC operates in high-dimensional binary space. SPPU operates
in low-dimensional balanced ternary space where the null state is a
first-class computational primitive.

\subsection{Ternary Spike-Based SNN (Neural Networks, 2025)}

\textbf{What:} Threshold-adaptive encoding converts analog to ternary
spikes. QT-SNN processes with quantized ternary weights. Speech and
EEG recognition.

\textbf{ArcLoom distinction:} Trained spiking neural network with
ternary encoding. Not structural perception. Not combinational.
Not training-free.

\subsection{Gap Analysis}

The following concepts return \textbf{zero results} in the literature:
\begin{itemize}[nosep]
  \item ``Coupling weights'' as a perception mechanism
  \item ``Dimensional exhaustion'' as a detection principle
  \item ``Geometric inevitability'' in hardware
  \item Perception that settles via coupling constraints without
        iteration or training
  \item $3^i$ positional weights as structural decoders in a
        perception fabric
\end{itemize}

ArcLoom's computational model --- combinational settling through
signed ternary coupling with dead-zone thresholding --- has
\nopriorart.

% ============================================================
\section{Dead Zone / Null State Computing}

\subsection{Three-Valued Logic}

Kleene and \L ukasiewicz three-valued logics are well-established in
formal logic. Hardware implementations exist (TDDFET gates, memristive
circuits, Huawei patent). In all cases, the third state represents
``unknown'' or ``high-impedance'' --- a failure to resolve.

\subsection{ArcLoom Distinction}

In the SPPU, the null trit is not ``unknown.'' It is ``the coupling
field has insufficient energy to collapse this dimension.'' This is
an \textit{active} computational state with specific consequences:

\begin{itemize}[nosep]
  \item It contributes zero to downstream coupling (multiplicative silence)
  \item It can be created by habituation (Krimelack dead-zone feedback)
  \item It can be created by insufficient input (sensor ambiguity)
  \item It is counted by L6 to detect dimensional exhaustion
  \item It can flip to committed ($+1$ or $-1$) when coupling pressure
        increases
\end{itemize}

No prior ternary architecture treats the null state as structural
uncertainty with these computational properties.

% ============================================================
\section{Summary: What ArcLoom Has That Nobody Else Has}

\begin{center}
\begin{longtable}{@{}p{4.5cm}p{4cm}p{5cm}@{}}
\toprule
\textbf{Capability} & \textbf{Nearest Prior Art} & \textbf{Gap} \\
\midrule
\endhead
Combinational ternary perception & Loihi 2 (async SNN) & No ML, no spikes, no clock \\
$3^i$ positional coupling weights & None & \nopriorart \\
Null trit as structural uncertainty & Three-valued logic & Not used as active computation \\
Krimelack structural motif memory & Memristor habituation & Not CAM, not pattern-level, not combinational \\
Folding division with native rounding & REBEL-2 (add/sub/mul only) & No division anywhere \\
L6 dimensional exhaustion & None & \nopriorart \\
Clockless habituation via dead-zone feedback & None & \nopriorart \\
Sensor-to-motor ternary chain & None & \nopriorart \\
Training-free embodied cognition & HDC (trains classifier) & End-to-end training-free \\
Autonomous obstacle navigation on ternary HW & None & \nopriorart \\
\bottomrule
\end{longtable}
\end{center}

% ============================================================
\section{Potential Allies / Collaborators}

\begin{center}
\begin{tabular}{@{}lll@{}}
\toprule
\textbf{Entity} & \textbf{What They Have} & \textbf{How They Help ArcLoom} \\
\midrule
USN Ternary Research Group & EDA tools (MRCS) & Synthesis and verification for tape-out \\
Prophesee & Event cameras (GenX320) & Async sensor input for SPPU \\
Efabless & Shuttle runs (SKY130, GF180) & ArcLoom-256 fabrication \\
Caltech/Yale async groups & ACT async design tools & Async BSIL timing verification \\
\bottomrule
\end{tabular}
\end{center}

% ============================================================
\section{Anticipated Objections and Responses}

\textbf{``Isn't this just another ternary computer?''}

No. The 5500FP is a ternary computer --- fetch, decode, execute.
ArcLoom has no instruction set. The answer settles through wire
propagation. Different computational model entirely.

\textbf{``How is this different from BitNet?''}

BitNet uses ternary values on binary hardware. ArcLoom uses ternary
logic as the native substrate. BitNet is a compression trick for
trained LLMs. ArcLoom is a perception architecture that requires
zero training.

\textbf{``Huawei already made a ternary chip.''}

Huawei patented ternary gates using standard ternary (0,1,2). ArcLoom
uses balanced ternary ($-1$,0,$+1$) which has the symmetric negation
property required for signed coupling. Huawei's chip (if it exists)
runs sequential instructions. ArcLoom settles combinationally.
Different number system, different architecture.

\textbf{``Neural accelerators already do low-power perception.''}

Neural accelerators require trained models, clocks, and sequential
matrix multiplication. SPPU requires no training, no clock, and no
multiplication in the decision path. Power is $10^{9}\times$ lower
than ARM Cortex-M0 for equivalent perception.

\textbf{``How do you know it works?''}

19,683 arithmetic operations verified on FPGA silicon. 41,430 division
tests at zero errors. Live sensor-to-motor chain demonstrated.
351-sample autonomous obstacle navigation with camera. Camera-sensor
ambiguity correctly perceived. See ASIC Spec v1.2, \S12
(Validation Status).

\textbf{``What about Loihi / Akida / SpiNNaker?''}

All three run trained spiking neural networks on async hardware. The
hardware is clockless; the computation is ML. SPPU is purely
combinational --- no spikes, no neurons, no weights learned from data.

% ============================================================
\section{DSF-AI as ML Model Auditor (New Capability, May 2026)}

During this competitive analysis, a new DSF-AI capability was
discovered: the UF-Core kernel (L0--L4) can be run directly on the
trained weight tensors of ML models. The kernel reads trained weights
as structural signals and extracts coupling profiles per tensor,
per layer, and per element.

\subsection{Demonstration: CathodeX MACE Ensemble Audit}

The DSF-AI kernel was run on the trained weights of the CathodeX
MACE-MP-0 ensemble (5 members, 5.5M parameters each) --- a materials
science ML model for cathode screening. Findings:

\begin{enumerate}[nosep]
  \item \textbf{Ensemble members are structurally identical.}
        All 5 members produce identical DSF profiles to 6 decimal
        places. The ensemble's epistemic uncertainty estimate
        (inter-model variance) may measure noise, not genuine
        disagreement.
  \item \textbf{Lithium has the weakest learned representation}
        of all cathode-relevant elements (coupling = 0.889,
        uncertainty = 0.434). The element most critical to cathode
        screening is the one the model knows least about.
  \item \textbf{Parameter budget allocated to wrong elements.}
        Top 5 by weight magnitude: S, As, Ta, Sb, Nb --- not
        cathode metals.
\end{enumerate}

\textbf{Implication:} DSF-AI can audit any ML model by reading its
weight structure. No other tool provides this capability. This is a
natural extension of the kernel's domain-agnostic structural field
extraction --- trained weights are just another signal.

See \textit{DSF-AI MACE Audit Report} (May 2026) for full details.

\vspace{2em}
\hrule
\vspace{0.5em}
{\color{tfe-gray}\small This document is confidential and for internal
use only. The competitive landscape analysis is based on publicly
available information as of May 2026. All ArcLoom claims reference
validated demonstrations documented in the ASIC Architecture
Specification v1.2. The ArcLoom SPPU architecture is the intellectual
property of Joseph Forrester / DSF-AI. The coupling weight derivation
methodology, UF pipeline internals, and coherence field physics are
trade secrets and are not disclosed in this document.}

\end{document}
